\chapter*{Introducción}
\addcontentsline{toc}{chapter}{Introducción}
\markboth{\MakeUppercase{Introducción}}{\MakeUppercase{Introducción}}

El auge del aprendizaje automático y en particular del aprendizaje profundo, ha transformado numerosas industrias y campos de investigación. Modelos como las redes neuronales artificiales han demostrado una capacidad sin precedentes para extraer patrones complejos de grandes volúmenes de datos, impulsando avances significativos en áreas como la visión por computadora, el procesamiento del lenguaje natural y la medicina. Su éxito radica en la habilidad de aprender directamente de los datos, ajustando millones de parámetros para realizar tareas de predicción, clasificación o regresión con una precisión notable.

Sin embargo, la creciente dependencia de estos modelos ha generado preocupaciones fundamentales en torno a la privacidad y la confidencialidad. El procesamiento de información sensible, como registros médicos, datos financieros o información personal, por parte de servicios de aprendizaje automático alojados en la nube o gestionados por terceros, expone a los usuarios a riesgos de fuga o uso indebido de sus datos. La necesidad de proteger la privacidad sin sacrificar la utilidad de los modelos de aprendizaje automático ha impulsado la investigación en técnicas de cómputo que preservan la privacidad.

En este contexto, el cifrado homomórfico emerge como una solución prometedora. El cifrado homomórfico permite realizar operaciones computacionales directamente sobre datos cifrados, de tal manera que el resultado de dichas operaciones, una vez descifrado, es idéntico al que se habría obtenido de realizar las mismas operaciones sobre los datos en texto plano. Esta propiedad única ofrece un mecanismo robusto para procesar información sensible en entornos no confiables, asegurando que los datos permanezcan cifrados durante todo el ciclo de procesamiento.

A pesar de su potencial, la aplicación práctica del cifrado homomórfico, especialmente en tareas computacionalmente intensivas como la inferencia de redes neuronales, ha estado limitada por su alto costo computacional. Los esquemas de cifrado homomórfico más avanzados, como CKKS (Cheon-Kim-Kim-Song), son capaces de manejar aritmética aproximada sobre números reales y complejos, lo cual los hace particularmente adecuados para las operaciones de punto flotante características del aprendizaje automático. No obstante, las implementaciones existentes de CKKS suelen depender de la CPU, lo que introduce cuellos de botella significativos cuando se trabaja con tensores grandes o modelos complejos.

El presente trabajo aborda esta limitación proponiendo el desarrollo de una biblioteca en Python que aproveche la capacidad de procesamiento paralelo de las Unidades de Procesamiento Gráfico (GPU) para acelerar las operaciones tensoriales bajo el esquema de cifrado homomórfico CKKS. Nuestro objetivo es facilitar la inferencia privada de modelos de aprendizaje automático, ofreciendo una herramienta eficiente que combine la seguridad criptográfica con el rendimiento computacional necesario para aplicaciones del mundo real. Para lograr esto se plantean los siguientes objetivos específicos:

\begin{itemize}
   \item Diseñar una interfaz en Python que permita interactuar con una biblioteca de cifrado homomórfico CKKS acelerada por GPU y emplear las operaciones básicas del esquema.
   \item Implementar las clases y módulos que permitan realizar operaciones con tensores y matrices, abstrayendo del usuario la complejidad del esquema CKKS.
   \item Crear casos de uso que ejemplifiquen el empleo de la biblioteca para aprovechar el cifrado homomórfico durante la inferencia en modelos de aprendizaje automático.
   \item Comparar el rendimiento de las redes neuronales implementadas con la biblioteca propuesta frente a implementaciones equivalentes construidas sobre otras bibliotecas de cifrado homomórfico.
\end{itemize}

El trabajo se aborda en dos fases. La primera es el desarrollo de la biblioteca, guiado por pruebas: se diseña primero la interfaz pública, se escriben las pruebas de las operaciones de mayor complejidad antes de implementarlas y se itera sobre la implementación hasta satisfacerlas. La segunda fase es la evaluación, para la cual se construye un banco de pruebas que ejecuta dos arquitecturas representativas, una red totalmente conectada aplicada a un problema de regresión y una red convolucional aplicada a uno de clasificación, sobre cuatro implementaciones distintas: la biblioteca propuesta, una ejecución de referencia en texto plano y dos bibliotecas de cifrado homomórfico del estado del arte. Todas las implementaciones parten de los mismos pesos entrenados y cada una se ejecuta en un proceso aislado, de manera que las diferencias observadas se atribuyan al esquema de cifrado y no al proceso de entrenamiento. La evaluación contempla la fidelidad de los resultados cifrados respecto de los obtenidos en texto plano, que mide la pérdida de precisión inducida por las aproximaciones polinómicas y por el ruido del esquema, y el rendimiento en tiempo de ejecución y consumo de memoria.

Tecnológicamente, la biblioteca se implementa en Python y delega el cómputo criptográfico en HEonGPU, una biblioteca escrita en C++ y CUDA que ejecuta las primitivas del esquema CKKS sobre la GPU y ofrece soporte de \textit{bootstrapping}. Los modelos se entrenan en texto plano con PyTorch, que sirve además como ejecución de referencia sin cifrado, y las bibliotecas de contraste son Orion, que emplea el mismo esquema CKKS pero sobre la CPU, y Concrete~ML, basada en el esquema TFHE.

El trabajo se organiza en los siguientes capítulos:

\begin{itemize}
   \item \textbf{Capítulo 1: Conceptos de Cifrado}:
Este capítulo introduce las nociones básicas del cifrado. Se presenta la definición formal de un esquema de cifrado, los dos grandes enfoques con que se construyen los esquemas modernos y la limitación común a todos los esquemas convencionales: la imposibilidad de operar sobre los datos cifrados sin descifrarlos previamente. Esta limitación motiva el cifrado totalmente homomórfico, sobre el que se articula este trabajo.

   \item \textbf{Capítulo 2: Cifrado Homomórfico}:
Este capítulo introduce los conceptos esenciales del cifrado homomórfico. Se exploran sus orígenes, la problemática del ruido inherente a estos esquemas, las diferencias entre cifrado parcialmente, algo y totalmente homomórfico, y las técnicas de gestión del ruido como el \textit{bootstrapping}. Se detallan los problemas matemáticos subyacentes, como el aprendizaje con errores y el aprendizaje con errores sobre anillos, y se presenta en profundidad el esquema CKKS, incluyendo sus mecanismos de codificación, generación de claves y operaciones homomórficas.

   \item \textbf{Capítulo 3: Cifrado Homomórfico sobre Redes Neuronales}:
Aquí se examina la intersección entre el cifrado homomórfico y el aprendizaje automático. Se discute cómo las arquitecturas de redes neuronales, tanto profundas como convolucionales, pueden adaptarse para operar sobre datos cifrados. Se profundiza en la implementación de operaciones lineales clave (producto punto, multiplicación matriz-vector y convoluciones) en el dominio homomórfico, y se aborda el desafío de las funciones de activación no lineales y sus aproximaciones polinómicas.

   \item \textbf{Capítulo 4: Marco Metodológico}:
Este aplica el marco metodológico: la infraestructura de hardware y software empleada, el diseño y la implementación de la biblioteca propuesta, las bibliotecas de referencia con las que se compara, los conjuntos de datos y las arquitecturas de redes neuronales sobre las que se evalúa, y la metodología de medición y evaluación.

   \item \textbf{Conclusiones y Recomendaciones}:
Finalmente se presentan las conclusiones y recomendaciones del trabajo. Se recogen los hallazgos derivados del diseño, la implementación y la evaluación de la biblioteca, se señalan las limitaciones del alcance abordado y se proponen las líneas de continuación identificadas.
\end{itemize}


